% Options for packages loaded elsewhere
% Options for packages loaded elsewhere
\PassOptionsToPackage{unicode}{hyperref}
\PassOptionsToPackage{hyphens}{url}
%
\documentclass[
  ignorenonframetext,
  aspectratio=169,
  russian,
  aspectratio=169]{beamer}
\newif\ifbibliography
\usepackage{pgfpages}
\setbeamertemplate{caption}[numbered]
\setbeamertemplate{caption label separator}{: }
\setbeamercolor{caption name}{fg=normal text.fg}
\beamertemplatenavigationsymbolshorizontal
% remove section numbering
\setbeamertemplate{part page}{
  \centering
  \begin{beamercolorbox}[sep=16pt,center]{part title}
    \usebeamerfont{part title}\insertpart\par
  \end{beamercolorbox}
}
\setbeamertemplate{section page}{
  \centering
  \begin{beamercolorbox}[sep=12pt,center]{section title}
    \usebeamerfont{section title}\insertsection\par
  \end{beamercolorbox}
}
\setbeamertemplate{subsection page}{
  \centering
  \begin{beamercolorbox}[sep=8pt,center]{subsection title}
    \usebeamerfont{subsection title}\insertsubsection\par
  \end{beamercolorbox}
}
% Prevent slide breaks in the middle of a paragraph
\widowpenalties 1 10000
\raggedbottom
\AtBeginPart{
  \frame{\partpage}
}
\AtBeginSection{
  \ifbibliography
  \else
    \frame{\sectionpage}
  \fi
}
\AtBeginSubsection{
  \frame{\subsectionpage}
}
\usepackage{iftex}
\ifPDFTeX
  \usepackage[T1]{fontenc}
  \usepackage[utf8]{inputenc}
  \usepackage{textcomp} % provide euro and other symbols
\else % if luatex or xetex
  \usepackage{unicode-math} % this also loads fontspec
  \defaultfontfeatures{Scale=MatchLowercase}
  \defaultfontfeatures[\rmfamily]{Ligatures=TeX,Scale=1}
\fi
\usepackage{lmodern}

\usefonttheme{serif} % use mainfont rather than sansfont for slide text
\ifPDFTeX\else
  % xetex/luatex font selection
  \setmainfont[]{DejaVu Serif}
  \setsansfont[]{DejaVu Sans}
  \setmonofont[]{DejaVu Sans Mono}
\fi
% Use upquote if available, for straight quotes in verbatim environments
\IfFileExists{upquote.sty}{\usepackage{upquote}}{}
\IfFileExists{microtype.sty}{% use microtype if available
  \usepackage[]{microtype}
  \UseMicrotypeSet[protrusion]{basicmath} % disable protrusion for tt fonts
}{}


\usepackage{longtable,booktabs,array}
\usepackage{calc} % for calculating minipage widths
\usepackage{caption}
% Make caption package work with longtable
\makeatletter
\def\fnum@table{\tablename~\thetable}
\makeatother
\usepackage{graphicx}
\makeatletter
\newsavebox\pandoc@box
\newcommand*\pandocbounded[1]{% scales image to fit in text height/width
  \sbox\pandoc@box{#1}%
  \Gscale@div\@tempa{\textheight}{\dimexpr\ht\pandoc@box+\dp\pandoc@box\relax}%
  \Gscale@div\@tempb{\linewidth}{\wd\pandoc@box}%
  \ifdim\@tempb\p@<\@tempa\p@\let\@tempa\@tempb\fi% select the smaller of both
  \ifdim\@tempa\p@<\p@\scalebox{\@tempa}{\usebox\pandoc@box}%
  \else\usebox{\pandoc@box}%
  \fi%
}
% Set default figure placement to htbp
\def\fps@figure{htbp}
\makeatother



\ifLuaTeX
\usepackage[bidi=basic,provide=*]{babel}
\else
\usepackage[bidi=default,provide=*]{babel}
\fi
\ifPDFTeX
\else
\babelfont{rm}[]{DejaVu Serif}
\fi
% get rid of language-specific shorthands (see #6817):
\let\LanguageShortHands\languageshorthands
\def\languageshorthands#1{}


\setlength{\emergencystretch}{3em} % prevent overfull lines

\providecommand{\tightlist}{%
  \setlength{\itemsep}{0pt}\setlength{\parskip}{0pt}}



 

\usepackage[]{csquotes}

\IfFileExists{plex-otf.sty}{
  %% Full TeXlive
  % \usepackage[%
  %   % math,
  %   RM={Scale=0.94},SS={Scale=0.94},SScon={Scale=0.94},TT={Scale=MatchLowercase,FakeStretch=0.9},DefaultFeatures={Ligatures=Common}
  % ]{plex-otf}
}{
  %% TinyTeX
  \usepackage{libertine}
}

%%% Load theme
% https://deic.uab.cat/~iblanes/beamer_gallery/
\IfFileExists{beamerthemegotham.sty}{
  %% Full TeXlive
  \usetheme{gotham}
  \gothamset{
    numbering=totalpagenumber,
    parttocframe default=off,
    sectiontocframe default=off,
    subsectiontocframe default=off,
  }
}{
  %% TinyTeX
  \usetheme{Madrid}
}
\makeatletter
\@ifpackageloaded{caption}{}{\usepackage{caption}}
\AtBeginDocument{%
\ifdefined\contentsname
  \renewcommand*\contentsname{Содержание}
\else
  \newcommand\contentsname{Содержание}
\fi
\ifdefined\listfigurename
  \renewcommand*\listfigurename{Список иллюстраций}
\else
  \newcommand\listfigurename{Список иллюстраций}
\fi
\ifdefined\listtablename
  \renewcommand*\listtablename{Список таблиц}
\else
  \newcommand\listtablename{Список таблиц}
\fi
\ifdefined\figurename
  \renewcommand*\figurename{Рисунок}
\else
  \newcommand\figurename{Рисунок}
\fi
\ifdefined\tablename
  \renewcommand*\tablename{Таблица}
\else
  \newcommand\tablename{Таблица}
\fi
}
\@ifpackageloaded{float}{}{\usepackage{float}}
\floatstyle{ruled}
\@ifundefined{c@chapter}{\newfloat{codelisting}{h}{lop}}{\newfloat{codelisting}{h}{lop}[chapter]}
\floatname{codelisting}{Список}
\newcommand*\listoflistings{\listof{codelisting}{Листинги}}
\makeatother
\makeatletter
\makeatother
\makeatletter
\@ifpackageloaded{caption}{}{\usepackage{caption}}
\@ifpackageloaded{subcaption}{}{\usepackage{subcaption}}
\makeatother

\usepackage{bookmark}
\IfFileExists{xurl.sty}{\usepackage{xurl}}{} % add URL line breaks if available
\urlstyle{same}
\hypersetup{
  pdftitle={Выполнение лабораторной работы №2},
  pdfauthor={Карпухин Клим},
  pdflang={ru-RU},
  hidelinks,
  pdfcreator={LaTeX via pandoc}}


\title{Выполнение лабораторной работы №2}
\subtitle{Первоначальная настройка git}
\author{Карпухин Клим}
\date{2026-03-01}

\begin{document}
\frame{\titlepage}


\section{Содержание}\label{ux441ux43eux434ux435ux440ux436ux430ux43dux438ux435}

\begin{frame}{Содержание}
\begin{enumerate}[<+->]
\tightlist
\item
  Информация о докладчике
\item
  Вводная часть и актуальность
\item
  Объект и предмет исследования
\item
  Цель работы
\item
  Задачи
\item
  Материалы, методы и инструменты
\item
  Ход работы (этапы, скриншоты)
\item
  Результаты и анализ
\item
  Выводы
\end{enumerate}
\end{frame}

\section{Информация}\label{ux438ux43dux444ux43eux440ux43cux430ux446ux438ux44f}

\begin{frame}{Докладчик}
\phantomsection\label{ux434ux43eux43aux43bux430ux434ux447ux438ux43a}
\begin{columns}[c]
\begin{column}{0.65\linewidth}
\begin{itemize}[<+->]
\tightlist
\item
  \textbf{Карпухин Клим}
\item
  Российский университет дружбы народов
\item
  Email:
  \href{mailto:1032255580@rudn.ru}{\nolinkurl{1032255580@rudn.ru}}
\item
  Роли: студент (лабораторная работа по ОС/виртуализации)
\end{itemize}
\end{column}

\begin{column}{0.35\linewidth}

\includegraphics[width=0.9\linewidth,height=\textheight,keepaspectratio]{./image/dokladchik.png}
\end{column}
\end{columns}
\end{frame}

\section{Вводная
часть}\label{ux432ux432ux43eux434ux43dux430ux44f-ux447ux430ux441ux442ux44c}

\begin{frame}{Актуальность}
\phantomsection\label{ux430ux43aux442ux443ux430ux43bux44cux43dux43eux441ux442ux44c}
\begin{itemize}[<+->]
\tightlist
\item
  Системы контроля версий (VCS) --- стандарт современной разработки.
\item
  Git является наиболее распространённой распределённой VCS.
\item
  Платформа GitHub --- ключевой инструмент для хостинга проектов и
  командной работы.
\item
  Владение Git необходимо для учёбы, исследований и будущей
  профессиональной деятельности.
\end{itemize}
\end{frame}

\begin{frame}{Объект и предмет исследования}
\phantomsection\label{ux43eux431ux44aux435ux43aux442-ux438-ux43fux440ux435ux434ux43cux435ux442-ux438ux441ux441ux43bux435ux434ux43eux432ux430ux43dux438ux44f}
\begin{itemize}[<+->]
\tightlist
\item
  \textbf{Объект:} система контроля версий Git и платформа GitHub.
\item
  \textbf{Предмет:} процесс первичной настройки окружения для работы с
  Git, включая аутентификацию и верификацию.
\end{itemize}
\end{frame}

\section{Цель и
задачи}\label{ux446ux435ux43bux44c-ux438-ux437ux430ux434ux430ux447ux438}

\begin{frame}{Цель}
\phantomsection\label{ux446ux435ux43bux44c}
Изучить идеологию и применение средств контроля версий, освоить базовые
навыки работы с git.
\end{frame}

\begin{frame}{Задачи}
\phantomsection\label{ux437ux430ux434ux430ux447ux438}
\begin{itemize}[<+->]
\tightlist
\item
  Создать базовую конфигурацию для работы с git.
\item
  Создать ключ SSH и ключ PGP.
\item
  Настроить автоматические подписи коммитов.
\item
  Зарегистрироваться на GitHub.
\item
  Создать локальный каталог для выполнения заданий по предмету.
\end{itemize}
\end{frame}

\section{Материалы и
методы}\label{ux43cux430ux442ux435ux440ux438ux430ux43bux44b-ux438-ux43cux435ux442ux43eux434ux44b}

\begin{frame}[fragile]{Материалы и методы}
\phantomsection\label{ux43cux430ux442ux435ux440ux438ux430ux43bux44b-ux438-ux43cux435ux442ux43eux434ux44b-1}
\begin{itemize}[<+->]
\tightlist
\item
  \textbf{Инструменты}

  \begin{itemize}[<+->]
  \tightlist
  \item
    Git --- распределённая система контроля версий.
  \item
    Утилита \texttt{gh} (GitHub CLI) для работы с GitHub.
  \item
    SSH для безопасного подключения к удалённым репозиториям.
  \item
    GnuPG (gpg) для создания ключей подписи коммитов.
  \end{itemize}
\item
  \textbf{Методы:} настройка конфигурационных файлов, генерация ключей
  шифрования, работа с командной строкой Linux.
\end{itemize}
\end{frame}

\section{Ход работы --- настройка Git
(этапы)}\label{ux445ux43eux434-ux440ux430ux431ux43eux442ux44b-ux43dux430ux441ux442ux440ux43eux439ux43aux430-git-ux44dux442ux430ux43fux44b}

\begin{frame}[fragile]{Этап 1: Установка программного обеспечения}
\phantomsection\label{ux44dux442ux430ux43f-1-ux443ux441ux442ux430ux43dux43eux432ux43aux430-ux43fux440ux43eux433ux440ux430ux43cux43cux43dux43eux433ux43e-ux43eux431ux435ux441ux43fux435ux447ux435ux43dux438ux44f}
\begin{itemize}[<+->]
\tightlist
\item
  Установлены \texttt{git} и \texttt{gh} с помощью пакетного менеджера
  \texttt{dnf}.
\end{itemize}

\begin{figure}

\centering{

\includegraphics[width=0.3\linewidth,height=\textheight,keepaspectratio]{image/1.png}

}

\caption{\label{fig-001}Установка git и gh}

\end{figure}%
\end{frame}

\begin{frame}{Этап 2: Базовая настройка Git}
\phantomsection\label{ux44dux442ux430ux43f-2-ux431ux430ux437ux43eux432ux430ux44f-ux43dux430ux441ux442ux440ux43eux439ux43aux430-git}
\begin{itemize}[<+->]
\tightlist
\item
  Заданы имя и email владельца репозитория.
\item
  Настроена кодировка UTF-8, имя начальной ветки, параметры autocrlf и
  safecrlf.
\end{itemize}

\begin{figure}

\centering{

\includegraphics[width=0.3\linewidth,height=\textheight,keepaspectratio]{image/2.png}

}

\caption{\label{fig-002}Базовая настройка Git}

\end{figure}%
\end{frame}

\begin{frame}[fragile]{Этап 3: Создание SSH-ключа}
\phantomsection\label{ux44dux442ux430ux43f-3-ux441ux43eux437ux434ux430ux43dux438ux435-ssh-ux43aux43bux44eux447ux430}
\begin{itemize}[<+->]
\tightlist
\item
  Сгенерированы ключи по алгоритмам \texttt{rsa} (4096 бит) и
  \texttt{ed25519}.
\end{itemize}

\begin{figure}

\centering{

\includegraphics[width=0.3\linewidth,height=\textheight,keepaspectratio]{image/3.png}

}

\caption{\label{fig-003}Создание SSH-ключей}

\end{figure}%
\end{frame}

\begin{frame}{Этап 4: Создание PGP-ключа}
\phantomsection\label{ux44dux442ux430ux43f-4-ux441ux43eux437ux434ux430ux43dux438ux435-pgp-ux43aux43bux44eux447ux430}
\begin{itemize}[<+->]
\tightlist
\item
  Создан ключ для подписи коммитов (RSA and RSA, 4096 бит).
\end{itemize}

\begin{figure}

\centering{

\includegraphics[width=0.3\linewidth,height=\textheight,keepaspectratio]{image/5.png}

}

\caption{\label{fig-005}Создание PGP-ключа}

\end{figure}%
\end{frame}

\begin{frame}{Этап 5: Добавление PGP-ключа в GitHub}
\phantomsection\label{ux44dux442ux430ux43f-5-ux434ux43eux431ux430ux432ux43bux435ux43dux438ux435-pgp-ux43aux43bux44eux447ux430-ux432-github}
\begin{itemize}[<+->]
\tightlist
\item
  Ключ скопирован и добавлен в веб-интерфейсе GitHub.
\end{itemize}

\begin{figure}

\centering{

\includegraphics[width=0.3\linewidth,height=\textheight,keepaspectratio]{image/9.png}

}

\caption{\label{fig-009}Добавление PGP-ключа в GitHub}

\end{figure}%
\end{frame}

\begin{frame}{Этап 6: Настройка автоматических подписей коммитов}
\phantomsection\label{ux44dux442ux430ux43f-6-ux43dux430ux441ux442ux440ux43eux439ux43aux430-ux430ux432ux442ux43eux43cux430ux442ux438ux447ux435ux441ux43aux438ux445-ux43fux43eux434ux43fux438ux441ux435ux439-ux43aux43eux43cux43cux438ux442ux43eux432}
\begin{itemize}[<+->]
\tightlist
\item
  Указан PGP-ключ для подписи, включена автоматическая подпись коммитов.
\end{itemize}

\begin{figure}

\centering{

\includegraphics[width=0.3\linewidth,height=\textheight,keepaspectratio]{image/10.png}

}

\caption{\label{fig-010}Настройка подписей коммитов}

\end{figure}%
\end{frame}

\begin{frame}[fragile]{Этап 7: Создание репозитория курса на основе
шаблона}
\phantomsection\label{ux44dux442ux430ux43f-7-ux441ux43eux437ux434ux430ux43dux438ux435-ux440ux435ux43fux43eux437ux438ux442ux43eux440ux438ux44f-ux43aux443ux440ux441ux430-ux43dux430-ux43eux441ux43dux43eux432ux435-ux448ux430ux431ux43bux43eux43dux430}
\begin{itemize}[<+->]
\tightlist
\item
  Использована команда \texttt{gh\ repo\ create} с шаблоном.
\end{itemize}

\begin{figure}

\centering{

\includegraphics[width=0.3\linewidth,height=\textheight,keepaspectratio]{image/12.png}

}

\caption{\label{fig-012}Создание репозитория из шаблона}

\end{figure}%
\end{frame}

\begin{frame}[fragile]{Этап 8: Настройка каталога курса и первый коммит}
\phantomsection\label{ux44dux442ux430ux43f-8-ux43dux430ux441ux442ux440ux43eux439ux43aux430-ux43aux430ux442ux430ux43bux43eux433ux430-ux43aux443ux440ux441ux430-ux438-ux43fux435ux440ux432ux44bux439-ux43aux43eux43cux43cux438ux442}
\begin{itemize}[<+->]
\tightlist
\item
  Создана структура каталогов, выполнены \texttt{git\ add},
  \texttt{commit}, \texttt{push}.
\end{itemize}

\begin{figure}

\centering{

\includegraphics[width=0.3\linewidth,height=\textheight,keepaspectratio]{image/14.png}

}

\caption{\label{fig-014}Первый коммит и push}

\end{figure}%
\end{frame}

\section{Результаты и
анализ}\label{ux440ux435ux437ux443ux43bux44cux442ux430ux442ux44b-ux438-ux430ux43dux430ux43bux438ux437}

\begin{frame}{Анализ и практическая значимость}
\phantomsection\label{ux430ux43dux430ux43bux438ux437-ux438-ux43fux440ux430ux43aux442ux438ux447ux435ux441ux43aux430ux44f-ux437ux43dux430ux447ux438ux43cux43eux441ux442ux44c}
\begin{itemize}[<+->]
\tightlist
\item
  \textbf{Результат:} полностью настроенное рабочее окружение для
  дальнейшей работы с курсом.
\item
  Созданы и добавлены в аккаунт GitHub ключи для безопасной
  аутентификации (SSH) и верификации авторства (PGP).
\item
  Все последующие коммиты из этого репозитория будут автоматически
  подписаны и помечены на GitHub как \enquote{Verified}.
\item
  \textbf{Практическая значимость:}

  \begin{itemize}[<+->]
  \tightlist
  \item
    Полученные навыки являются базой для выполнения всех последующих
    лабораторных работ.
  \item
    Обеспечена безопасность и прослеживаемость изменений в проектах.
  \end{itemize}
\end{itemize}
\end{frame}

\section{Выводы}\label{ux432ux44bux432ux43eux434ux44b}

\begin{frame}{Выводы}
\phantomsection\label{ux432ux44bux432ux43eux434ux44b-1}
\begin{itemize}[<+->]
\tightlist
\item
  В ходе работы была изучена идеология систем контроля версий.
\item
  Освоены практические навыки первичной настройки Git.
\item
  Сгенерированы SSH- и PGP-ключи, необходимые для безопасной работы и
  подтверждения авторства.
\item
  Настроена интеграция локального Git-клиента с платформой GitHub.
\item
  Создано рабочее пространство с применением шаблона репозитория,
  готовое для выполнения заданий по предмету.
\item
  \textbf{Цель лабораторной работы достигнута полностью.}
\end{itemize}
\end{frame}




\end{document}
